\section{OpenAI's construction for a design}\label{sec:setup}

We recall the construction of \cite{OpenAI2026} in the generality in which we use it, and we prove for an
arbitrary design the degree bound, the bounds for the denominators at the finite places and the
nonvanishing at the primes $N=p$. Carreira \cite{Carreira2026} checks that the argument of OpenAI goes
through for designs satisfying a few conditions; we need the same for the finite places and, in addition,
for the entries modulo $p^2$.

\subsection{The determinant}
Let $D=(n_0,b_0,q_0,g_0,h_0)$ be a design, with the notation of Section~\ref{sec:overview}. For $0\le r<n$
the row polynomials are
\begin{equation}\label{eq:rows}
 P_r(t)=(1-t)^ht^{C-1}T_{|r-g|}(1/t),\qquad \deg P_r<H,
\end{equation}
together with companion polynomials $D_r$ built from $U_{|r-g|-1}$ \cite[(3)]{OpenAI2026}. Thus
$P_r$ depends only on $|r-g|$. Put $f(t)=\sqrt{1-t^2}$ and $d\mu(t)=|t|\,dt/f(t)$ on $(-1,1)$, and
\[
 A_r(t)=P_r(t)-\frac32\,\mathbf 1_{\{t>0\}}\frac{f(t)}{t}D_r(t),\qquad
 \psi_k(s)=s^{b+k}(1-s)^q .
\]
The determinant studied by OpenAI is that of the $n\times n$ matrix
\begin{equation}\label{eq:matrix}
 \cA_N(r,k)=\int_{-1}^1\int_0^1 A_r(t)\,\frac{\psi_k(s)}{1-ts}\,ds\,d\mu(t)
 \qquad(0\le r,k<n).
\end{equation}
The factor $(1-s)^q$ in $\psi_k$ amounts to a finite difference of order $q$ applied to the columns.
The entries are rational linear combinations of $1$ and $G$: there are rational matrices $\cA^0_N$
and $\cL_N$ with
\begin{equation}\label{eq:L}
 \cA_N=\cA^0_N+G\,\cL_N,\qquad
 \cL_N(r,k)=\sum_{v=0}^{q}(-1)^v\binom qv\sum_{i}[t^i]P_r\cdot4\,c_{(i-b-k-v)/2},
\end{equation}
where $c_l=4^{-l}\binom{2l}{l}$ for integers $l\ge0$ and $c_z=0$ otherwise
\cite[Prop.~2.1 and (20)]{OpenAI2026}. The rational part is
\begin{equation}\label{eq:rational-part}
 \begin{gathered}
 \cA^0_N(r,k)=\sum_{v=0}^q(-1)^v\binom qv R(r,b+k+v),\\
 R(r,j)=\sum_{i<H}P_r[i]\,M(i,j)-\frac32\sum_{i<H}D_r[i]\,Z(i,j)
 \end{gathered}
\end{equation}
with rational numbers $M(i,j)$ (moments of $\mu$ against the Cauchy kernel) and $Z(i,j)$, where
$Z(i,i)=-\sum_{k\le i}k^{-2}$ and $Z(i,j)=(H_i-H_j)/(i-j)$ for $i\ne j$, $H_i=\sum_{k\le i}1/k$
\cite[Prop.~2.1]{OpenAI2026}. The design enters \eqref{eq:rational-part} only through the integer
coefficients $P_r[i],D_r[i]$ and the range of $j$, which is $[b,L)$. Put
$\Delta_N(X)=\det(\cA^0_N+X\cL_N)\in\Q[X]$, so that $\Delta_N(G)=\det\cA_N$.

\subsection{The degree}
\begin{lemma}\label{lem:degree}
$\deg\Delta_N\le n-g$.
\end{lemma}

\begin{proof}
By \eqref{eq:L} the row $r$ of $\cL_N$ depends only on $P_r$, and hence only on $|r-g|$. So the rows
$r$ and $2g-r$ of $\cL_N$ agree for $0\le r<g$, and $2g-r\le n-1$ because $2g<n$. Subtracting row $r$ from
row $2g-r$ for each $r<g$ does not change the determinant and makes $g$ rows of the new coefficient of
$X$ vanish. A determinant $\det(B+XL')$ in which $L'$ has $n-g$ nonzero rows has degree at most $n-g$ in $X$.
\end{proof}

\subsection{Rank deficiency and valuation}
We use repeatedly the following elementary fact.

\begin{lemma}\label{lem:rank-valuation}
Let $M$ be an $n\times n$ matrix with entries in $\Z_{(p)}$ whose reduction modulo $p$ has rank at most
$n-e$. Then $v_p(\det M)\ge e$.
\end{lemma}

\begin{proof}
Factor the reduction through $\F_p^{n-e}$ and lift the two factors to integral matrices $U,V$, so that
$M=UV+pE$ with $E$ integral at $p$. Expanding the determinant multilinearly in the columns, every term
that uses more than $n-e$ columns of $UV$ vanishes, and every other term has at least $e$ factors $p$.
\end{proof}

\subsection{The digit formula}
The reduction of $p^2\cA^0_N$ modulo $p$ is described by the digits in base $p$ of the coefficients of the
row polynomials. This is the one place where the arithmetic of the moments enters; it is design-free.
For an odd prime $p$ and $c\in\F_p^n$ put
\[
 P_c=\sum_rc_rP_r,\qquad D_c=\sum_rc_rD_r,\qquad W=(1-X^2)^{(p-1)/2},
\]
(reading the $P_r,D_r$ as polynomials in $X$), and let $F[k]$ denote the coefficient of $X^k$.

\begin{lemma}[digit formula]\label{lem:digits}
Let $p$ be an odd prime with $H<p^2$, let $0\le j<H$, and write $j=Jp+\ell$ with $0\le\ell<p$. Then
$p^2R(r,j)$ is $p$-integral for every row $r$, and for every $c\in\Z^n$
\[
 p^2\sum_rc_rR(r,j)\equiv\sum_{u\ge0}\Bigl(\alpha_{u,J}\,(XWP_c)[up+\ell]+\beta_{u,J}\,D_c[up+\ell]\Bigr)
 \pmod p ,
\]
where $\alpha_{u,J}=M(u-1,J)$ and $\beta_{u,J}=-\tfrac32Z(u,J)$ are $p$-integral numbers that do not depend on
$c$, on $\ell$ or on the design.
\end{lemma}

\begin{proof}
Two facts about the numbers $M$ and $Z$ are needed, for $0\le i,j<p^2$. Write $j=Jp+\ell$ and let
$\varrho\in[0,p)$ be the residue of $j-i-1$ modulo $p$, $u=(i+1+\varrho-\ell)/p$. First, $p^2M(i,j)$ is
$p$-integral and $p^2M(i,j)\equiv W[\varrho]\,M(u-1,J)\pmod p$; this is the digit reduction of the moments in
\cite[\S3]{OpenAI2026}, obtained from the $p$-adic analysis of their recursions, and it does not involve a design.
Second, $p^2Z(i,j)$ is $p$-integral and is congruent to $Z(\lfloor i/p\rfloor,J)$ if $i\equiv j\pmod p$ and to
$0$ otherwise. For the second, write $i'=\lfloor i/p\rfloor$. Since $i<p^2$, $H_i=p^{-1}H_{i'}+S_i$ and
$H^{(2)}_i=p^{-2}H^{(2)}_{i'}+S'_i$ with $S_i,S'_i$ $p$-integral, which gives the case $i=j$. If $i\not\equiv j$,
then $p^2Z(i,j)=p\cdot pZ(i,j)\equiv0$. If $i\equiv j$, $i\ne j$, then $i-j=(i'-J)p$, and the integers
$k\in(j,i]$ prime to $p$ form $i'-J$ complete systems of residues modulo $p$, so that
$\sum_k1/k\equiv(i'-J)\sum_{r=1}^{p-1}1/r\equiv0\pmod p$; hence $H_i-H_j\equiv(H_{i'}-H_J)/p$ and
$p^2Z(i,j)\equiv Z(i',J)$.

Now use \eqref{eq:rational-part}. By the first fact the part $\sum_iF[i]\,p^2M(i,j)$ with $F=P_c$ is
$\sum_iF[i]\,W[\varrho(i)]\,M(u(i)-1,J)$. For fixed $u$ the indices with $u(i)=u$ are those for which
$\varrho=up+\ell-1-i$ lies in $[0,p-1]$, and $W$ has degree $p-1$, so the sum over these $i$ is
$(XFW)[up+\ell]$; this gives the $\alpha$-terms. By the second fact the part with $D_c$ is
$\sum_{i\equiv\ell}D_c[i]\,Z(\lfloor i/p\rfloor,J)=\sum_uZ(u,J)D_c[up+\ell]$, which gives the $\beta$-terms.
The $p$-integrality of $M(u-1,J)$ and $Z(u,J)$ for $u,J<p$ is part of the first and second facts.
\end{proof}

For the filtered columns we use the following form. Let $0\le k<n$, $s=b+k$, $J=\lfloor s/p\rfloor$ and
$\mu=(s\bmod p)+q$. If $\mu<p$, the $q+1$ raw columns $s,\dots,s+q$ that enter column $k$ have the same
digit $J$ and residues $(s\bmod p)+v<p$. Since
$\sum_{v=0}^q(-1)^v\binom qvF[t+v]=[X^{t+q}]\bigl((X-1)^qF\bigr)$ for every polynomial $F$,
Lemma~\ref{lem:digits} gives
\begin{equation}\label{eq:filtered-digits}
 p^2\sum_rc_r\cA^0_N(r,k)\equiv\sum_{u\ge0}\Bigl(\alpha_{u,J}\,Q_c[up+\mu]+\beta_{u,J}\,\widehat D_c[up+\mu]\Bigr),
\end{equation}
where $Q_c=(X-1)^q\,XWP_c$ and $\widehat D_c=(X-1)^qD_c$. A combination $c$ whose responses vanish on all
raw columns in $[b,L)$ lies in the left kernel of $\cA^0_{N,p}$, and a combination whose responses on the
columns outside a set vanish and which kills the set by \eqref{eq:filtered-digits} does as well.

\subsection{Integrality and the valuation at $p$}
\begin{proof}[Proof of Proposition~\ref{prop:valuation}]
By Lemma~\ref{lem:digits} the entries of $p^2\cA^0_N$ are $p$-integral, since the filter is an integer
combination of raw columns with $j<L\le H$. The entries of $\cL_N$ are $p$-integral for odd $p$, since
$P_r[i]\in\Z$ and $c_l\in\Z[1/2]$, so $p^2\cL_N\equiv0\pmod{p^2}$ and, for $z$ with denominator prime to
$p$, $p^2(\cA^0_N+z\cL_N)$ has $p$-integral entries with the same reduction as $p^2\cA^0_N$. By
Lemma~\ref{lem:rank-valuation} the determinant of the latter has valuation at least $\nu_N(p)$, and
$\det\bigl(p^2(\cA^0_N+z\cL_N)\bigr)=p^{2n}\Delta_N(z)$.
\end{proof}

\subsection{The denominators}
Let $p$ be an odd prime with $H<p^2$ and $p\le H$. The method of \cite[\S3]{OpenAI2026}, applied to a
design, bounds from below the valuation of $\Delta_N(z)$ for every $z$ with denominator prime to $p$ by
$-\ell_N(p)$, where $\ell_N(p)$ counts the columns whose entries have a pole at $p$, once the coefficients of
the row polynomials are grouped by their digits in base $p$ and the column filter is taken into account. It
is given by two formulas, in which all differences are truncated at $0$. For $2p\le H$,
\[
 \ell_N(p)=2n-\Bigl(\min\bigl(p,\,L-p,\,A\bigr)-\max\bigl(b,\,H-2p\bigr)-q\Bigr),
\]
and for $H<2p\le2H$, with $R=(L-b)-(\min(A,L-p)-b)-(\min(p,L)-\max(b,H-p))$,
$P'=\min(A,L-p)-b$ and $E'=\min(b,H-p)-(L-p)$,
\[
 \ell_N(p)=\min\bigl(2n,\ n+R,\ 2R+P'+E'\bigr).
\]
These are the exponents of the generic valuation theorems in the formalization; for the primes that
interest us they are the same as in OpenAI's paper for its design. For the design $D$ the function
$\ell_N(p)=Nf(p/N)$ is piecewise linear in $x=p/N$ ($\lfloor\cdot\rfloor$ rounding aside):
\[
 f(x)=\begin{cases}
 116&x\le59/2,\\
 175-2x&59/2\le x\le69/2,\\
 106&69/2\le x\le39,\\
 184-2x&39\le x\le81/2,\\
 265-4x&81/2\le x\le48,\\
 217-3x&48\le x\le69,\\
 148-2x&69\le x\le70,\\
 78-x&70\le x\le78,
 \end{cases}
\]
and $f=0$ for $x\ge78$, since for $p>H$ the matrix $\cA^0_N+z\cL_N$ is $p$-integral. Its integral over
$[1,78]$ is $24269/4$, and $f=116=2\cdot58$ on $[0,1]$. Note that $f(x)=2n_0$ for $x\le59/2$, which
contains all the primes in $(N,2N]$ and, in particular, the whole range $[12,26]$ of the rank
deficiencies of Theorem~\ref{thm:content}; there the bound of Proposition~\ref{prop:valuation} improves
$-\ell_N(p)=-2n$ to $\nu_N(p)-2n$.

The primes $p\le\sqrt H$ are not covered by these bounds. For them we use the cruder bound
$6n\log_pH$ for the exponent, which contributes $O(\sqrt N\,N\log N)=o(N^2)$ in total.

\subsection{The prime $2$}
At $p=2$ the entries are not integral and the bound comes from the expansion of the determinant into
minors, in the manner of \cite[\S3]{OpenAI2026}. Its exponent is governed by a quadratic form in the
numbers $m,\ell$ of columns of two exceptional types, which is
\[
 \tfrac32m^2+\tfrac12\ell^2-(H+C-b)\,m-C\,\ell+Cn ,
\]
and the following identity shows that it is at least $-\mathsf A_2$ whenever $m,\ell\ge0$,
$m+\ell\le n$ and $C\ge n$, where
\begin{equation}\label{eq:A2}
 \mathsf A_2=\frac{(n+H-b)^2}8-\frac{n^2}2 .
\end{equation}
With $s=n-m-l$ and $m_*=(n+H-b)/4$,
\[
 \tfrac32m^2+\tfrac12\ell^2-(H+C-b)m-C\ell+Cn+\mathsf A_2
 =2(m-m_*)^2+s\,(C-n+m)+\tfrac12s^2\ge0 .
\]
This gives a $2$-adic bound $v_2(\Delta_N(i))\ge-\mathsf A_2-O(N\log N)$, uniform in the integer $i$. For $D$,
$\mathsf A_2=(58+78-9)^2/8-58^2/2=2673/8$ in units of $N^2$.

\subsection{The common denominator}
\begin{proof}[Proof of Proposition~\ref{prop:denominators}]
Let $\kappa_N(p)$ be the largest of the bounds of Theorem~\ref{thm:content} that apply at $p$ (and $0$ if
none does), for primes $p\le H$ with $H<p^2$, and put $e_N(p)=\min\bigl(\ell_N(p),\,2n-\kappa_N(p)\bigr)$. For the
primes with $p^2\le H$ use the crude bound above, for $p=2$ the $2$-adic bound, and let $D_N$ be the
product of $p^{e_N(p)}$ over $p\le H$. By Proposition~\ref{prop:valuation} and the bounds of this section,
$v_p(\Delta_N(i))\ge-e_N(p)$ for every integer $i\ge0$ and every $p$ (for $p>H$ the matrix is
$p$-integral), so $D_N\Delta_N(i)\in\Z$. A prime $q\in(N,2N]$ has $\kappa_N(q)=0$, since no window
contains it, and $\ell_N(q)=2n$, so $q^{2n}\mid D_N$ (here $N\ge79$ guarantees $q>\sqrt H$).

For the size of $D_N$, the sum of $\ell_N(p)\log p$ over the odd primes is, by the prime number theorem in
the weighted form
\begin{equation}\label{eq:weighted-pnt}
 \frac1N\sum_{p}F(p/N)\log p\longrightarrow\int_0^\infty F(x)\,dx\qquad(N\to\infty),
\end{equation}
for bounded, Riemann integrable, compactly supported $F$, equal to $N^2\int_0^{78}f(x)\,dx+o(N^2)$ with
$\int_0^{78}f=116+24269/4$. The primes with $\kappa_N(p)>0$ lie in $(12N,26N]$ and have
$\ell_N(p)=116N$, so their exponent is $116N-\kappa_N(p)$; the sum of $\kappa_N(p)\log p$ is
$N^2\int\kappa+o(N^2)$ by \eqref{eq:weighted-pnt}. The $2$-adic part contributes
$\mathsf A_2\log2\,N^2+o(N^2)$.
\end{proof}

The polynomial $\Delta_N$ has degree at most $n-g=d_0N$ and takes values in $D_N^{-1}\Z$ at the
integers. An integer-valued polynomial of degree $m$ has coefficients in $\frac1{m!}\Z$ (Newton's
interpolation formula), so
\begin{equation}\label{eq:QN}
 Q_N=(d_0N)!\,D_N\,\Delta_N(X)\in\Z[X],\qquad \deg Q_N\le d_0N,
\end{equation}
and $\log(d_0N)!=o(N^2)$. In this way the gain from the rank deficiencies enters through the integer
polynomial itself, and not through a separate common factor of its coefficients.

\begin{lemma}\label{lem:recycling}
For all large $N$ and every prime $N<q\le2N$, the $j$-th coefficient of $Q_N$ is divisible by $q^{2j}$.
\end{lemma}

\begin{proof}
For these primes the matrices $q^2\cA^0_N$ and $\cL_N$ are $q$-integral by Lemma~\ref{lem:digits}, and
$q^{2n}$ divides $D_N$. The identity
\[
 Q_N(X/q^2)=(d_0N)!\,\frac{D_N}{q^{2n}}\det\bigl(q^2\cA^0_N+X\cL_N\bigr)
\]
shows that $Q_N(X/q^2)$ has $q$-integral coefficients; its $j$-th coefficient is $c_{N,j}/q^{2j}$.
\end{proof}

\subsection{Nonvanishing}
\begin{proposition}\label{prop:nonvanishing}
Let $D$ be a design and suppose that the three $n_0\times n_0$ matrices $B_\sigma$, $\sigma\in\{0,1,-1\}$,
defined below, are invertible. Then for all large primes $p$ and every rational $z$ such that $pz$ has
denominator prime to $p$, $\Delta_p(z)\ne0$. In particular $\Delta_p(a/b)\ne0$ whenever $a\in\Z$, $b\ge1$
and $p^2\nmid b$.
\end{proposition}

The matrices are those of OpenAI's reduction modulo $p$ at $N=p$, with the design in place of $48$: let
$B$ be the rational $(n_0+1)\times n_0$ matrix of the entries of the filtered rational part at $N=1$ for
the rows $0,\dots,n_0$ (the row $n_0$ is an auxiliary row), and let $B_\sigma$ have rows
$B_r+\sigma B_{r+1}$, $0\le r<n_0$. For the design $D$, $B$ has $59$ rows and $58$ columns.

\begin{proof}
Let $p>4H_0$, with $H_0=H/N=78$, so that $H=H_0p<p^2/4$, and let $p$ be prime to the numerators and denominators
of the entries of $B$ and of the three determinants. In the construction at $N=p$ write the row index
$r=pr_0+i$ with $0\le r_0<n_0$ and $0\le i<p$. Since $q=q_0p$, $C=C_0p$, $h=h_0p$, $g=g_0p$, the Frobenius
relation $(1-t)^{p}\equiv1-t^p$ expresses the row polynomials modulo $p$ through those of the design at
$N=1$ and the successor row $r_0+1$; this is where the auxiliary row enters, and it is valid for rows on
both sides of $g$ because the polynomials are defined through $T_{r-g}$ and $U_{r-g-1}$ for signed
indices. The filter satisfies $s^{b_0p+k}(1-s)^{q_0p}=s^{\ell}(s^p)^{b_0+k_0}(1-s^p)^{q_0}$ in $\F_p[s]$, so
the filtered columns of one residue class use the raw columns $b_0+k_0,\dots,b_0+k_0+q_0$ of the design at
$N=1$ and never change the residue. By \cite[Lemma~3.1 and eqs.~(31), (35), (36)]{OpenAI2026}, which
need only $H<p^2/4$, the reduction of $p^2\cA^0_p$ modulo $p$ is then, after ordering the rows by $(i,r_0)$ and
the columns by $(\ell,k_0)$, the matrix $a^{\mathsf T}\otimes B_0+b^{\mathsf T}\otimes B_1$, where $a,b$ are the
fixed $p\times p$ matrices of OpenAI's identity (51), with $\det a=2^{-p(p-1)/2}$. The permutation
matrix $\Pi$ in this decomposition has the eigenvalue $0$ once and $\pm1$ each $(p-1)/2$ times, and
therefore
\[
 \det=(\det a)^{n_0}\,\det B_0\;\det(B_0+B_1)^{(p-1)/2}\det(B_0-B_1)^{(p-1)/2}\ne0
\]
in $\F_p$. Hence $p^{2n}\Delta_p(0)$ is a $p$-adic unit. For $z$ with $pz$ of denominator prime to $p$, the
difference $p^2z\cL_p=p\,(pz)\,\cL_p$ is $p$-integral with zero reduction (the entries of $\cL_p$ are
$p$-integral), so the reduction of $p^2(\cA^0_p+z\cL_p)$ is the same invertible matrix and
$v_p(\Delta_p(z))=-2n$. If $p^2\nmid b$, then $p(a/b)$ has denominator prime to $p$.
\end{proof}

For the design $D$ the invertibility of $B_0$, $B_0+B_1$ and $B_0-B_1$ is certified by exhibiting their
inverses modulo the prime $1009$, which exceeds $H_0=78$ so that every denominator of $B$ is a unit
there; the three determinants are then nonzero rational numbers. The computation modulo $1009$ certifies
the base matrices; the conclusion for all large primes is a theorem and not an extrapolation from
finitely many primes. Nothing in the proof uses the proportions of the design except $2g<n$, and $H<p^2/4$.
